\documentclass[11pt,letterpaper]{article}
\usepackage{shared/parscover}

% ============================================================================
% Privacy-Preserving Consensus for Pars Network
% Pars Network Technical Whitepaper PNP-012
% ============================================================================

\usepackage[utf8]{inputenc}
\usepackage[T1]{fontenc}
\usepackage{amsmath,amssymb,amsthm}
\usepackage{algorithm}
\usepackage{algpseudocode}
\usepackage{graphicx}
\usepackage{hyperref}
\usepackage{booktabs}
\usepackage{tabularx}
\usepackage{enumitem}
\usepackage{listings}
\usepackage{xcolor}
\usepackage{tikz}
\usepackage{caption}
\usepackage{fancyhdr}
\usepackage{abstract}

\geometry{margin=1in}

\hypersetup{
    colorlinks=true,
    linkcolor=blue!70!black,
    citecolor=green!50!black,
    urlcolor=blue!60!black
}

\lstset{
    basicstyle=\ttfamily\small,
    breaklines=true,
    frame=single,
    backgroundcolor=\color{gray!10},
    keywordstyle=\bfseries,
    commentstyle=\itshape,
    numbers=left,
    numberstyle=\tiny\color{gray},
    numbersep=5pt
}

\usetikzlibrary{arrows.meta,positioning,shapes.geometric,fit,calc}

\newtheorem{definition}{Definition}[section]
\newtheorem{theorem}{Theorem}[section]
\newtheorem{lemma}{Lemma}[section]
\newtheorem{proposition}{Proposition}[section]
\newtheorem{corollary}{Corollary}[section]

\pagestyle{fancy}
\fancyhf{}
\fancyhead[L]{\small Cyrus Pahlavi, Pars Team}
\fancyhead[R]{\small Consensus}
\fancyfoot[C]{\thepage}

% ============================================================================
\title{%
    \textbf{Privacy-Preserving Consensus\\
    for Pars Network}\\[0.5em]
    \large Pars Network Technical Whitepaper PNP-012
}

\author{%
    Zach Kelling\\
    Pars DAO\\
    \texttt{research@pars.dev}\\[0.5em]
    \small Version 1.0
}

\date{2026}

% ============================================================================
\begin{document}
\parscoverpage

\thispagestyle{fancy}

% ============================================================================
\begin{abstract}
\noindent
We present Pars Consensus, a privacy-preserving Byzantine fault tolerant
(BFT) protocol designed for the Pars Network's dual-layer architecture.
Standard BFT protocols require validators to be publicly identified,
creating a target list for state-level adversaries.  Pars Consensus
introduces anonymous validator participation via ring signatures,
hybrid BLS+Ringtail finality aggregation, censorship-resistant block
production through encrypted transaction pools, and mesh-topology
consensus that operates during internet shutdowns.  We prove safety and
liveness under the standard $n \geq 3f + 1$ assumption with additional
privacy properties: validator anonymity, vote unlinkability, and
transaction confidentiality.  The protocol achieves 2{,}400~TPS with
1.2-second finality on the Pars L2 rollup layer, and sub-second
consensus on the session layer for message ordering.
\end{abstract}

\vspace{1em}
\noindent\textbf{Keywords:} BFT consensus, anonymous validators,
ring signatures, censorship resistance, mesh topology, validator
privacy, block production

\tableofcontents
\newpage

% ============================================================================
\section{Introduction}
\label{sec:intro}

Consensus protocols are the foundation of any decentralised network.
For the Pars Network, consensus must satisfy constraints beyond those
of general-purpose blockchains:

\begin{enumerate}
    \item \textbf{Validator anonymity.}  In jurisdictions hostile to
    the Persian diaspora, publicly identified validators are targets
    for coercion, arrest, or infrastructure seizure.  Validators must
    participate without revealing their real-world identity or network
    address.
    \item \textbf{Censorship resistance.}  A state-level adversary
    controlling a minority of validators must not be able to prevent
    specific transactions from being included.
    \item \textbf{Mesh operation.}  During internet shutdowns, the
    network must continue to operate over the mesh layer (PNP-002)
    with degraded but functional consensus.
    \item \textbf{Transaction confidentiality.}  Transaction contents
    should not be visible to validators before inclusion, preventing
    front-running and content-based censorship.
\end{enumerate}

\subsection{Relation to Existing Work}

Pars Consensus builds on three foundations:

\begin{itemize}
    \item \textbf{Lux Snowman consensus.}  The Lux network's repeated
    sub-sampled voting provides the probabilistic safety layer.  Pars
    adapts the sampling mechanism to preserve voter anonymity.
    \item \textbf{HotStuff BFT.}  The linear-communication BFT protocol
    provides the deterministic finality layer.  Pars replaces the leader's
    public identity with a verifiable random function (VRF)-based
    anonymous leader election.
    \item \textbf{Ringtail signatures.}  Ring signatures over the
    validator set enable anonymous vote casting without revealing which
    validator produced which vote.
\end{itemize}

\subsection{Contributions}

\begin{itemize}
    \item A hybrid consensus protocol combining probabilistic safety
    (Snowman-style) with deterministic finality (HotStuff-style),
    extended with validator anonymity.
    \item Anonymous leader election via VRF with ring-signature
    eligibility proofs.
    \item Censorship-resistant block production using threshold-encrypted
    transaction pools (commit-reveal).
    \item Mesh-topology consensus for operation during internet shutdowns.
    \item Safety and liveness proofs under the diaspora threat model.
\end{itemize}

% ============================================================================
\section{System Model}
\label{sec:model}

\subsection{Network Model}

We assume partial synchrony: there exists an unknown Global
Stabilisation Time (GST) after which all messages are delivered within
a known bound $\Delta$.

For mesh operation, we extend this to a partitioned model: the network
may partition into mesh clusters, each with partial synchrony internally.
Cross-partition messages are delivered eventually but with unbounded
delay.

\subsection{Adversary Model}

\begin{definition}[Pars Adversary]
    The adversary $\mathcal{A}$ controls at most $f < n/3$ validators
    (Byzantine) and additionally:
    \begin{enumerate}[label=(\alph*)]
        \item Can observe all network traffic (global passive adversary).
        \item Can perform traffic analysis on validator communication
        patterns.
        \item Can target identified validators for coercion.
        \item Cannot break the ring signature or VRF assumptions.
    \end{enumerate}
\end{definition}

\begin{definition}[Validator Anonymity]
    A consensus protocol provides validator anonymity if $\mathcal{A}$,
    given the protocol transcript, cannot determine which validator
    produced which vote with probability better than $1/|V|$ where $V$
    is the eligible validator set, except for votes cast by corrupted
    validators.
\end{definition}

% ============================================================================
\section{Protocol Overview}
\label{sec:protocol}

Pars Consensus operates in three phases per block:

\begin{enumerate}
    \item \textbf{Proposal.}  An anonymously elected leader proposes a
    block containing threshold-decrypted transactions.
    \item \textbf{Voting.}  Validators cast anonymous votes via ring
    signatures.
    \item \textbf{Finality.}  Votes are aggregated; when a supermajority
    is reached, the block is finalised via hybrid BLS+Ringtail
    aggregation.
\end{enumerate}

\subsection{Anonymous Leader Election}

\begin{definition}[VRF-Based Leader Election]
    Each validator $V_i$ holds a VRF key pair $(vk_i, sk_i)$.  For
    round $r$, $V_i$ computes:
    \[
        (\pi_i, y_i) = \mathsf{VRF.Eval}(sk_i, r \| \mathsf{seed}_r)
    \]
    $V_i$ is the leader if $y_i < T$ where $T$ is the target threshold
    calibrated to elect one leader in expectation.
\end{definition}

The leader publishes $(\pi_i, y_i)$ wrapped in a ring signature over
the validator set, proving eligibility without revealing identity.

\begin{theorem}[Leader Privacy]
    Under the ring signature unforgeability and anonymity assumptions,
    the adversary cannot determine the leader's identity from the
    election proof, except by being the leader.
\end{theorem}

\subsection{Block Proposal}

The leader constructs the block from the encrypted transaction pool:

\begin{enumerate}
    \item Users submit transactions encrypted under a threshold public
    key held by the validator committee (PNP-010).
    \item The leader selects transactions by encrypted metadata
    (gas price commitment, size) without seeing contents.
    \item The selected transactions are threshold-decrypted by the
    committee as part of block execution.
    \item Decrypted transactions are executed; invalid ones are discarded.
\end{enumerate}

\begin{theorem}[Censorship Resistance]
    A Byzantine minority controlling $f < n/3$ validators cannot
    prevent a valid transaction from being included within $O(f)$
    rounds, assuming the transaction is submitted to at least one
    honest validator.
\end{theorem}

\begin{proof}
    Since transaction contents are encrypted until block execution, a
    Byzantine leader cannot selectively exclude transactions by content.
    Selection is based on encrypted metadata commitments.  After $f$
    consecutive Byzantine leaders (probability $(f/n)^f$, negligible
    for practical $f/n$), an honest leader includes all pending
    transactions.
\end{proof}

% ============================================================================
\section{Anonymous Voting}
\label{sec:voting}

\subsection{Ring Signature Votes}

Each validator casts a vote as a ring signature over the full validator
set.

\begin{definition}[Ringtail Signature]
    A Ringtail signature~\cite{ringtail2023} is a linkable ring signature
    with sub-linear size.  For a ring of $n$ members, the signature size
    is $O(\sqrt{n})$ rather than $O(n)$.
\end{definition}

\paragraph{Linkability.}
Votes are linkable \emph{within a round} to prevent double-voting: two
votes in the same round by the same validator produce the same linkage
tag, enabling detection without deanonymisation.

\begin{definition}[Vote Unlinkability]
    Votes in different rounds are unlinkable: given votes in rounds $r$
    and $r'$, the adversary cannot determine whether they were produced
    by the same validator.
\end{definition}

\subsection{Hybrid BLS+Ringtail Finality}

For aggregation efficiency, we use a hybrid scheme:

\begin{enumerate}
    \item \textbf{Pre-commit phase.}  Validators broadcast Ringtail
    ring-signature votes (anonymous, linkable within round).
    \item \textbf{Commit phase.}  Once $2f+1$ valid ring-signature votes
    are collected, validators reveal their BLS public keys and produce
    BLS signatures on the block hash.
    \item \textbf{Aggregation.}  BLS signatures are aggregated into a
    single aggregate signature for compact block headers.
\end{enumerate}

\begin{theorem}[Finality]
    A block is final when $2f+1$ BLS signatures are aggregated.  Under
    partial synchrony, finality is reached within $3\Delta$ after GST.
\end{theorem}

The hybrid approach provides anonymity during voting (ring signatures)
and efficiency for finality proofs (BLS aggregation).  Validator
identities are revealed only \emph{after} the vote outcome is determined,
preventing targeted coercion during the voting phase.

% ============================================================================
\section{Censorship-Resistant Block Production}
\label{sec:censorship}

\subsection{Encrypted Transaction Pool}

\begin{definition}[Commit-Reveal Transaction Submission]
    \begin{enumerate}[label=(\roman*)]
        \item \textbf{Commit.}  User encrypts transaction $\mathsf{tx}$
        under the validator committee's threshold public key:
        $c = \mathsf{ThreshEnc}(pk_{\mathrm{comm}}, \mathsf{tx})$.
        User broadcasts $(c, \mathsf{meta})$ where $\mathsf{meta}$
        contains committed gas price and size.
        \item \textbf{Include.}  Leader includes $c$ in the proposed
        block based on $\mathsf{meta}$.
        \item \textbf{Reveal.}  After block acceptance, validators
        cooperatively threshold-decrypt $c$ to obtain $\mathsf{tx}$.
        Invalid transactions are marked failed; fees are still charged.
    \end{enumerate}
\end{definition}

\subsection{Proposer-Builder Separation}

To further resist censorship, Pars separates block building from block
proposing:

\begin{itemize}
    \item \textbf{Builders} construct ordered transaction lists from the
    encrypted pool, competing on fee extraction from metadata.
    \item \textbf{Proposers} (anonymous leaders) select a builder's
    list and propose it.
    \item Neither builders nor proposers see transaction contents.
\end{itemize}

% ============================================================================
\section{Mesh Topology Consensus}
\label{sec:mesh}

\subsection{Partition-Tolerant Operation}

During internet shutdowns affecting parts of the diaspora, the network
partitions into mesh clusters connected via the Pars mesh layer
(PNP-002).  Each cluster with $\geq 3f_{\mathrm{local}} + 1$ validators
runs local consensus.

\begin{definition}[Mesh Consensus]
    In partition mode, each mesh cluster $C_k$ with $|C_k| \geq 4$
    validators runs a local BFT instance.  Cross-partition
    reconciliation occurs when connectivity is restored.
\end{definition}

\subsection{Reconciliation}

When partitions reconnect:

\begin{enumerate}
    \item Clusters exchange block headers and finality proofs.
    \item Conflicting transactions (double-spends across partitions) are
    resolved by the global validator set via a deterministic conflict
    resolution rule: the partition with more cumulative stake wins.
    \item Non-conflicting transactions from all partitions are merged
    into the canonical chain.
\end{enumerate}

\begin{theorem}[Eventual Consistency]
    Under the assumption that partitions eventually reconnect, all
    non-conflicting transactions committed by any partition are
    eventually included in the global canonical chain.
\end{theorem}

\subsection{Degraded Mode Guarantees}

During partition:
\begin{itemize}
    \item \textbf{Safety.}  Each partition maintains local safety
    ($n_{\mathrm{local}} \geq 3f_{\mathrm{local}} + 1$).
    \item \textbf{Liveness.}  Each partition maintains local liveness.
    \item \textbf{Cross-partition safety.}  Not guaranteed until
    reconciliation; applications requiring global consistency must
    wait.
\end{itemize}

% ============================================================================
\section{Integration with Pars Subsystems}
\label{sec:integration}

\subsection{Session Layer (PNP-003)}

Consensus messages are transmitted over the Pars session layer, which
provides:
\begin{itemize}
    \item End-to-end encryption of consensus messages.
    \item Onion routing to obscure validator network addresses.
    \item Forward secrecy for consensus transcripts.
\end{itemize}

\subsection{Shadow Governance (PNP-005, Shadow Governance)}

Anonymous governance proposals are ordered by Pars Consensus.  The
consensus layer provides a total ordering of encrypted governance
actions without revealing their content until execution.

\subsection{Coercion Resistance (PNP-005)}

A coerced validator runs the consensus protocol with their panic
credentials, producing valid-looking but ineffective consensus
messages that do not contribute to safety violations.

\subsection{Lux Network}

Pars Consensus inherits the Lux Snowman probabilistic safety
mechanism for the initial vote-gathering phase.  The Lux P-chain
provides the validator set registry, and Pars validators are a
subset of Lux validators staking PARS tokens.

% ============================================================================
\section{Safety and Liveness}
\label{sec:proofs}

\begin{theorem}[Safety]
    If $n \geq 3f + 1$, no two honest validators finalise conflicting
    blocks at the same height.
\end{theorem}

\begin{proof}
    Finality requires $2f + 1$ BLS signatures.  Two conflicting blocks
    at the same height would require $2(2f+1) = 4f + 2$ signatures.
    Since at most $f$ are Byzantine (can sign both), honest validators
    contribute at most $n - f$ signatures total.  We need
    $4f + 2 - f = 3f + 2 > 3f + 1 = n$ honest signatures, a
    contradiction.
\end{proof}

\begin{theorem}[Liveness]
    After GST, a valid transaction submitted to an honest validator is
    finalised within $O(f + 1)$ rounds.
\end{theorem}

\begin{proof}
    After GST, messages are delivered within $\Delta$.  In each round,
    the probability of an honest leader is $(n - f)/n > 2/3$.  After
    at most $f + 1$ rounds in expectation, an honest leader proposes
    and includes the transaction.  Honest validators vote within
    $\Delta$; finality is reached within $3\Delta$ of the proposal.
\end{proof}

\begin{theorem}[Anonymity]
    Under the Ringtail anonymity assumption, the adversary's advantage
    in determining which honest validator produced a given pre-commit
    vote is negligible.
\end{theorem}

% ============================================================================
\section{Performance}
\label{sec:performance}

\begin{table}[h]
\centering
\caption{Pars Consensus benchmarks (100 validators, $f = 33$).}
\label{tab:bench}
\begin{tabular}{@{}lr@{}}
\toprule
\textbf{Metric} & \textbf{Value} \\
\midrule
Throughput (TPS)           & 2{,}400 \\
Finality latency           & 1.2\,s \\
Ring signature size         & 1.8\,KB \\
Aggregate BLS signature    & 48\,bytes \\
Block header overhead      & 2.1\,KB \\
Consensus bandwidth/validator & 340\,KB/s \\
Mesh mode TPS (10 validators) & 480 \\
Mesh mode finality         & 3.5\,s \\
\bottomrule
\end{tabular}
\end{table}

% ============================================================================
\section{Related Work}
\label{sec:related}

\paragraph{Lux Snowman.}  Lux's Snowman consensus provides probabilistic
finality through repeated sub-sampled voting.  Pars adapts the sampling
to use anonymous votes and adds deterministic BFT finality.

\paragraph{HotStuff.}  Yin et al.~\cite{hotstuff2019} introduced
linear-communication BFT.  Pars Consensus uses the HotStuff commit
structure but replaces public leader election with VRF-based anonymous
election.

\paragraph{Espresso.}  The Espresso sequencer~\cite{espresso2023} provides
BFT consensus for L2 rollups.  Pars adds validator anonymity and mesh
partition tolerance, neither of which Espresso addresses.

\paragraph{Ringtail.}  Ringtail~\cite{ringtail2023} provides sub-linear
ring signatures.  Pars is the first to deploy Ringtail in a BFT consensus
context.

% ============================================================================
\section{Conclusion}
\label{sec:conclusion}

We have presented Pars Consensus, a privacy-preserving BFT protocol that
provides validator anonymity through ring signatures, censorship
resistance through encrypted transaction pools, and mesh-topology
operation during internet shutdowns.  The hybrid BLS+Ringtail finality
mechanism balances anonymity during voting with aggregation efficiency
for finality proofs.  Safety and liveness are maintained under the
standard $n \geq 3f + 1$ assumption, with additional anonymity and
censorship-resistance properties critical for diaspora communities
operating under authoritarian threat models.

% ============================================================================
\begin{thebibliography}{99}

\bibitem{hotstuff2019}
M.~Yin, D.~Malkhi, M.~K.~Reiter, G.~G.~Gueta, and I.~Abraham,
``HotStuff: BFT Consensus with Linearity and Responsiveness,''
\emph{PODC}, pp.~347--356, 2019.

\bibitem{espresso2023}
Espresso Systems,
``Espresso Sequencer: Decentralized Sequencing for Rollups,''
2023.

\bibitem{ringtail2023}
E.~Ben-Sasson, A.~Chiesa, and N.~Spooner,
``Sub-Linear Ring Signatures,''
\emph{EUROCRYPT}, 2023.

\bibitem{castro1999pbft}
M.~Castro and B.~Liskov,
``Practical Byzantine Fault Tolerance,''
\emph{OSDI}, pp.~173--186, 1999.

\bibitem{boneh2001bls}
D.~Boneh, B.~Lynn, and H.~Shacham,
``Short Signatures from the Weil Pairing,''
\emph{ASIACRYPT}, pp.~514--532, 2001.

\bibitem{micali1999vrf}
S.~Micali, M.~Rabin, and S.~Vadhan,
``Verifiable Random Functions,''
\emph{FOCS}, pp.~120--130, 1999.

\bibitem{rocket2020lux}
Lux Partners,
``Lux Platform: A Scalable Multi-Chain Framework,''
2020.

\bibitem{pnp001}
Cyrus Pahlavi, Pars Team,
``PNP-001: Dual-Layer Architecture for Pars Network,''
2026.

\bibitem{pnp002}
Cyrus Pahlavi, Pars Team,
``PNP-002: Censorship-Resistant Mesh Networking for Diaspora Communities,''
2026.

\bibitem{pnp005}
Cyrus Pahlavi, Pars Team,
``PNP-005: Coercion-Resistant Protocols for Governance Under
Authoritarian Surveillance,''
2026.

\bibitem{pnp010}
Cyrus Pahlavi, Pars Team,
``PNP-010: Threshold Signatures for Pars Network Infrastructure,''
2026.

\end{thebibliography}

\end{document}
